%% fuer Drucker "bilby":
%% dvips -O 5mm,1mm

%% fuer Drucker "goofy":
%% dvips600 -O 0mm,10mm Dagstuhl

\documentstyle[mathrsfs,colordvi,boldmath,semcolor,semlayer,portrait,rotate,german,epsf,fancybox,mathsym
]{seminar}
\input{unilogo}
\input{amssym.def}
\input{amssym}

\def\dummy#1{}
\overlaystrue
\def\overlay#1{{}}

\slidewidth9.5in
\slideheight6.7in
%\def\slidetopmargin{\iflandscape.2in\else0.8in\fi}
\def\slidetopmargin{\iflandscape.3in\else0.8in\fi}
\def\slidebottommargin{\iflandscape1.0in\else0.4in\fi}
\def\slideleftmargin{\iflandscape.2in\else.5in\fi}
\def\sliderightmargin{\iflandscape1.0in\else.8in\fi}

\slidesmag{4}
\centerslidesfalse
\slideframewidth0.5pt
%\slideframe{none}

%\portraitonly
%\landscapeonly
%\def\leftsliderotation#1{\rotate[l]{#1}}
%\sliderotation{left}
\rotateheaderstrue

%%%%%%%%%%%%%%%%%%%%%%%%%%%%%%%%%%%%%%%%%%%%%%
%% Makros
%%%%%%%%%%%%%%%%%%%%%%%%%%%%%%%%%%%%%%%%%%%%%%
\def\bra#1{\left<#1\right|}
\def\ket#1{\left|#1\right>}
\def\openone{\leavevmode\hbox{\small1\kern-3.3pt\normalsize1}}

\def\mathsymfont#1{\ifmmode{\mathchoice{%
    \mbox{\the\textfont0\ignorespaces#1}}{%
    \mbox{\the\textfont0\ignorespaces#1}}{%
    \mbox{\the\scriptfont0\ignorespaces#1}}{%
    \mbox{\the\scriptscriptfont0\ignorespaces#1}}}%
  \else{{\rm\ignorespaces#1}}\fi}
%\def\C{{\cal C}}
\def\Cdual{\C^{\perp}}
\def\Czdual{\C_2^{\perp}}
\def\Csum{\sum_{\bm{c}\in\C}}
\def\Cdualsum{\sum_{\bm{c}\in\Cdual}}
\def\F{\mathsymfont{I\kern-.25em F}}
\def\ds{\displaystyle}

\def\entspricht{\mathrel{\hat{=}}}

\def\si{id}
\def\sx{{\sigma_x}}
\def\sy{{\sigma_y}}
\def\sz{{\sigma_z}}

\def\w{\omega}
\def\ww{\overline{\omega}}

\def\dmin{\mathop{\rm d}_{\rm H min}\nolimits}
\def\dH{\mathop{\rm d}_{\rm H}\nolimits}
\def\dHnor{\mathop{\underline{\rm d}}_{\rm H}\nolimits}
\def\wH{\mathop{\rm w}_{\rm H}\nolimits}

\def\lowE{\raisebox{-0.9ex}{$\scriptstyle\sim$}\kern-0.6em E}

\def\trace{\mathop{\rm Tr}\nolimits}

\newfont{\smallbsf}{cmssbx10}
\newfont{\bsf}{cmssbx10 scaled 1400}
\newfont{\lbsf}{cmssbx10 scaled 1800}
\newfont{\emsf}{cmssi10}
\newfont{\trm}{cmr10}  % \ten\roman
\newfont{\twrm}{cmr12} % \twelve\roman
\newfont{\brm}{cmbx10 scaled 1400} % \bold\roman
\def\R{\mbox{\trm I\kern-.2em R}}
%\def\bm#1{\mbox{\boldmath$#1$}}
\newcommand{\heading}[1]{%
  \begin{center}
    \bsf
    \shadowbox{#1}%
  \end{center}
  \vspace{1ex minus 1ex}%
}
\newcommand{\leftheading}[1]{%
{    \bsf
    \shadowbox{#1}%
}
  \vspace{1ex minus 1ex}%
}
\newbox{\logobox}
\setbox\logobox=\hbox{\raise 0.17\baselineskip\hbox{\epsfbox{logo.eps}}}
\def\logo{\usebox{\logobox}}
\newpagestyle{MH}{\hss\hbox to \hsize{PhysComp96\hss\thepage}\hss}%
{\hss\hbox to \hsize{\logo Graduiertenkolleg 12/1999}
P.~Wocjan\hss}

\newpagestyle{nohead}{}{\hss\hbox to \hsize{\tiny\logo
{IAKS} \thepage/PhysComp96}\hss}

\newpagestyle{nohead_beth}{}{}

%{\hss\hbox to
%  \hsize{\tiny\ \logo\ IAKS\hfil P.~Wocjan}\hss}

\newbox\abc
\newdimen\xxx
\long\def\tbox#1{\leavevmode\setbox\abc\hbox{#1}\xxx\ht\abc\dp\abc0pt\ht\abc0pt\raise-\xxx\box\abc}

%%%%%%%%%%%%%%%%%%%%%%%%%%%%%%%%%%%%%%%%%%%%%%

\begin{document}
\frenchspacing
\sf

%%\pagestyle{nohead}
\pagestyle{nohead_beth}
%%\pagestyle{empty}

\begin{slide*}\sf
\begin{center}{\bsf Suboptimal Decoding}\end{center}
\vfill
For all block codes used on a memoryless classical channel the optimal decision
rule is the \emph{Maximum-likelihood} decoding. In general, the optimal 
decision rule on classical-quantum channels is not known.

Holevo has introduced the following suboptimal decoding rule 
${\mathscr D}_{{\mathcal H}}=
({\mathbf D}_1,{\mathbf D}_2,\ldots,{\mathbf D}_M)$ for any block code 
${\mathscr B}$ used on c-q channels ${\mathscr C}$ with arbitrary signal 
states $S_i$ . The decision operators are given by
$$
   {\mathbf D}_m =
     \left(\sum_{j=1}^M {\mathbf S}_j^r\right)^{-{1\over 2}}
     {\mathbf S}_m^r \,
     \left(\sum_{j=1}^M {\mathbf S}_j^r\right)^{-{1\over 2}}\,,
$$
$m=1,2\ldots,M$, and $r\in (0,1]$. 
This decision rule produces a word error probability $P_{\rm em}$ 
upperbounded by
$$
  P_{\rm em} \le \tilde{P}_{\rm em} =
    \sum_{j=1\atop j\neq m}^M
      {\rm Tr}\, {\mathbf S}_m^{1-r} {\mathbf S}_j^r \,.
$$
\vfill
\end{slide*}


\begin{slide*}\sf
\begin{center}{\bsf Union Bound}\end{center}
\vfill
{\bf Lemma 1}
The word error probability $P_{\rm em}$ of the binary
code $[N,R]$ used on the binary classical-quantum channel ${\mathscr C}$ and 
decoded by the suboptimal decision rule ${\mathscr D}_{\mathcal H}$ with 
$r=\frac{1}{2}$ is upperbounded by
$$ P_{\rm em} \le \sum_{j=1\atop j\neq m}^M 
     2^{-N [-\dHnor({\mathbf x}_m,{\mathbf x}_j) \log_2 (c)]}\,, $$
where the channel parameter $c:=\trace \sqrt{S_1} \sqrt{S_2}$ and 
$\dHnor({\mathbf x}_m,{\mathbf x}_j)=\dH({\mathbf x}_m,{\mathbf x}_j)/N$ is the 
normalized Hamming distance between the code words ${\mathbf x}_m$ and 
${\mathbf x}_j$.
\vfill
Consequently, the overall block decoding error probability is upperbounded 
by the \emph{union bound} given by
$$ P_{\rm e} \le \tilde{P}_{\rm e} = 
   \frac{1}{M} \sum_{m=1}^M \sum_{j=1\atop j\neq m}^M 
           2^{-N [-\dHnor({\mathbf x}_m,{\mathbf x}_j)\log_2 (c)]}\,. $$
\vfill
\end{slide*}

\begin{slide*}\sf
\begin{center}{\bsf Rescaled Binomial Average Multiplicity Enumerators}
\end{center}
\vfill
{\bf Lemma 2}
If the average multiplicity enumerators of the codes in the fixed rate
sequence ${\rm FRS}(R)$ are given by
rescaled binomial coefficients
$$
  \left\{ \underline{M}_n^{(i)} = {a_i(n)\over 2^{N_i (1-R)}}  
  {N_i\choose n}\right\}_{n=0}^{N_i}\,,\quad 
  \sum_{n=0}^{N_i} \underline{M}_n^{(i)} = M_i
$$
where $a_i(n)=o(2^n)$, than the corresponding IAAME function is given by
$$
  {\mathscr M}(\underline{\delta}_{\rm H},R)_{\rm FRS}=
    {\rm H}_2(\underline{\delta}_{\rm H}) - (1 - R) \,,
$$
where ${\rm H}_2(x)=-x\log_2(x) - (1-x)\log_2 (1-x)$ represents the binary 
entropy function.
\vfill
\end{slide*}

\begin{slide*}\sf
\begin{center}{\bsf Cutoff Rate Lower Bound}\end{center}
\vfill
{\bf Theorem 1} (Cutoff rate lower bound)
Fixed rate sequences with rescaled binomial average multiplicity enumerators 
used on the binary c-q channel ${\mathscr C}$ and decoded by the suboptimal 
rule ${\mathscr D}_{\mathcal H}$ with $r=\frac{1}{2}$ have an error exponent 
whose lower bound is the cutoff rate bound given by 
$$  \lowE_{\rm cut}(R) = R_0 - R $$
where $R_0$ is the cutoff rate given by
$$ R_0 = 1-\log_2(1+c)\,. $$
\vfill
\emph{Only} code words on the effective Hamming distance
$$ \underline{\delta}_{\rm H eff}={c\over 1+c} $$
determine the cutoff rate bound. Note that this distance does not depend on 
the code rate.
\vfill
\end{slide*}

\begin{slide*}\sf
\begin{center}{\bsf Gilbert-Varshamov Bound}\end{center}
\vfill
The \emph{Gilbert-Varshamov lower bound} on the asymptotic 
minimal normalized distance of binary block codes $[N,R]$ 
is given by
$$ \underline{\delta}_{\rm GV}(R) = {\rm H}_2^{-1}(1-R)\,, $$
where ${\rm H}_2^{-1}$ denotes the inverse of the binary entropy 
function in the interval $[0,{1\over 2}]$. 

The IAAME of a FSR($R$) with rescaled binomial average multiplicity enumerators 
satisfying the Gilbert-Varshamov bound is given by the 
{\red \emph{expurgated IAAME}}
$$
  {\mathscr M}(\underline{\delta}_{\rm H},R)_{\rm FSR} = \left\{ 
  \begin{array}{cl}
    -\infty & \mbox{for } \underline{\delta}_{\rm H} \le 
    \underline{\delta}_{\rm GV}\,, \\
    {\rm H}_2(\underline{\delta}_{\rm H}) - (1 - R) & \mbox{otherwise}.
   \end{array} \right.
$$
\vfill
{\bf Lemma 3}
Fixed rate sequences of linear block codes with rescaled binomial multiplicity
enumerators satisfy the Gilbert-Varshamov bound.
\vfill
\end{slide*}

\begin{slide*}\sf
\begin{center}{\bsf Expurgated Lower Bound}\end{center}
\vfill
{\bf Theorem 2} (Expurgated lower bound)
Fixed rate sequences with rescaled binomial multiplicity enumerators 
satisfying the Gilbert-Varshamov bound used on the binary c-q channel 
${\mathscr C}$ and decoded by the suboptimal rule ${\mathscr D}_{{\mathcal H}}$ 
with $r=\frac{1}{2}$ have an error exponent whose lower bound is the 
expurgated bound given by
$$ \lowE_{\rm ex}(R)=\log_2\left({1\over c}\right) \cdot 
  \underline{\delta}_{\rm GV}(R) $$
for rates $R$ below the expurgated rate $R_{\rm ex}$
$$ R_{\rm ex} = 1-{\rm H}_2\left(\underline{\delta}_{\rm H eff}\right)\, $$
and the cutoff rate bound for rates above $R_{\rm ex}$.
\vfill
\end{slide*}

\begin{slide*}\sf
\begin{center}{\bsf Conclusion}\end{center}
\vfill
A conceptually simple method for derivation of lower bounds on the error 
exponent of specific families of block codes used on classical-quantum 
channels with arbitrary signal states over a finite Hilbert space has been 
presented.
It has benn shown that families of binary block codes with appropriately 
rescaled binomial multiplicity enumerators used on binary
c-q channels and decoded by the suboptimal decision rule 
introduced by Holevo attain the expurgated and cutoff rate lower bounds on the 
error exponent.
\vfill
For the classical binary symmetric channel is known that code
families with rescaled binomial multiplicity enumerators attain the tight
part of the random coding bound and thus the channel capacity. This has been 
obtained by an improved upperbounding of the decoding error probability 
assuming the Maximum-likelihood decoding. Unfortunately, this method cannot be 
applied directly to binary c-q channels since the optimal
decision rule for these channels is not known.
\vfill
\end{slide*}

\end{document}









